\documentclass[12pt,a4paper]{article}
\usepackage[utf8]{inputenc}
\usepackage[T1]{fontenc}
\usepackage[polish]{babel}
\usepackage{geometry}
\usepackage{enumitem}
\usepackage{titlesec}
\usepackage{parskip}
\usepackage{fancyhdr}
\usepackage{xcolor}
\usepackage{mdframed}
\usepackage{microtype}
\usepackage{booktabs}
\usepackage{hyperref}
\usepackage{multirow}
\usepackage{array}
\usepackage{longtable}

\geometry{margin=2.5cm}

\hypersetup{colorlinks=true, linkcolor=black, urlcolor=blue}

\pagestyle{fancy}
\fancyhf{}
\rhead{SDLC-LLM Benchmark}
\lhead{Zadanie R3 -- Conflict Detection}
\rfoot{\thepage}

\titleformat{\section}{\large\bfseries}{}{0em}{}[\titlerule]
\titleformat{\subsection}{\normalsize\bfseries}{}{0em}{}
\titleformat{\subsubsection}{\normalsize\itshape}{}{0em}{}

\definecolor{metabg}{gray}{0.93}
\definecolor{promptbg}{RGB}{235,245,255}
\definecolor{gptcolor}{RGB}{16,163,127}
\definecolor{claudecolor}{RGB}{210,105,30}
\definecolor{llamacolor}{RGB}{100,100,180}
\definecolor{score3}{RGB}{220,240,220}
\definecolor{score2}{RGB}{255,243,205}
\definecolor{score1}{RGB}{255,220,200}
\definecolor{score0}{RGB}{255,200,200}

\newcommand{\scorebox}[1]{%
  \ifnum#1=3\colorbox{score3}{\textbf{3}}%
  \else\ifnum#1=2\colorbox{score2}{\textbf{2}}%
  \else\ifnum#1=1\colorbox{score1}{\textbf{1}}%
  \else\colorbox{score0}{\textbf{0}}\fi\fi\fi
}

\begin{document}

% ─── STRONA TYTUŁOWA ────────────────────────────────────────────────
\begin{center}
  {\LARGE\bfseries Requirements Artifact}\\[1.5em]
  {\large \textbf{Case study:} Appointment Booking System}\\[0.3em]
  {\large \textbf{Model:} gemini-3-flash-preview \quad }\\[0.3em]
  {\large \textbf{Tryb:} LLM-only \quad }\\[0.3em]

\end{center}

\vspace{1em}\hrule\vspace{1.5em}

% ─── PROMPT ─────────────────────────────────────────────────────────
\section{Prompt}

\begin{mdframed}[backgroundcolor=metabg, linewidth=0pt, innerleftmargin=10pt, innerrightmargin=10pt, innertopmargin=8pt, innerbottommargin=8pt]
\textbf{System prompt (simplified):}\\
\small Jesteś doświadczonym analitykiem wymagań. Odpowiadaj WYŁĄCZNIE czystym, w pełni poprawnym kodem JSON. Nie dodawaj żadnych znaczników Markdown.

[KROK 1] 
Na podstawie opisu systemu wygeneruj role i wymagania funkcjonalne.

[KROK 2] 
Na podstawie opisu i wygenerowanych FR wygeneruj NFR i user stories.

[KROK 3] 
Na podstawie user stories wygeneruj acceptance criteria w formacie Given/When/Then.

\end{mdframed}


\section{Cel systemu}
System umożliwia użytkownikom rezerwację wizyt u specjalistów (np. lekarzy lub konsultantów). Użytkownicy mogą przeglądać dostępne terminy, dokonywać rezerwacji oraz je anulować. Specjaliści zarządzają swoim grafikiem. System powinien zapobiegać konfliktom rezerwacji i obsługiwać różne role użytkowników.

\section{Role}
\begin{imtemize}
    \item \textbf{Użytkownik} - przeglądanie\_terminów, rezerwacja\_wizyty, anulowanie\_własnej\_rezerwacji
    \item \textbf{Specjalista} - zarządzanie\_grafikiem, podgląd\_rezerwacji, anulowanie\_wizyty\_pacjenta
    \item \textbf{Administrator} - zarządzanie\_użytkownikami, konfiguracja\_systemu, raportowanie
\end{imtemize}


\section{Wymagania funkcjonalne}
\begin{description}[leftmargin=1.8cm, labelwidth=1.3cm, labelsep=0.5cm]
    \item[FR1] System musi umożliwiać użytkownikom wyszukiwanie i przeglądanie dostępnych terminów u wybranych specjalistów.
    \item[FR2] System musi umożliwiać użytkownikom dokonanie rezerwacji wybranego wolnego terminu.
    \item[FR3] System musi umożliwiać użytkownikom anulowanie istniejącej rezerwacji zgodnie z polityką anulowań.
    \item[FR4] System musi umożliwiać specjalistom definiowanie, edytowanie i usuwanie dostępnych slotów czasowych w ich grafiku.
    \item[FR5] System musi automatycznie blokować termin w momencie potwierdzenia rezerwacji, uniemożliwiając podwójne zapisy na tę samą godzinę.
    \item[FR6] System musi wymuszać autoryzację i weryfikację uprawnień dla każdej operacji zgodnie z przypisaną rolą użytkownika.
\end{description}

\section{Wymagania niefunkcjonalne}
\begin{description}[leftmargin=1.8cm, labelwidth=1.3cm, labelsep=0.5cm]
    \item[NFR1] System musi zapewniać czas odpowiedzi na zapytanie o dostępność terminów nie przekraczający 2 sekund przy 1000 równoczesnych użytkownikach.
    \item[NFR2] Wszystkie dane osobowe oraz dane o wizytach muszą być szyfrowane zgodnie ze standardem AES-256 oraz spełniać wymogi RODO.
    \item[NFR3] System musi charakteryzować się dostępnością na poziomie 99,9\% w skali miesiąca.
    \item[NFR4] Interfejs użytkownika musi być w pełni responsywny (RWD), zapewniając poprawną obsługę na urządzeniach mobilnych i desktopowych.
    \item[NFR5] System musi automatycznie wylogować użytkownika po 15 minutach nieaktywności w celu zapewnienia bezpieczeństwa sesji.
\end{description}

\section{User stories}
\begin{description}[leftmargin=1.8cm, labelwidth=1.3cm, labelsep=0.5cm]
    \item[US1] Jako Użytkownik chcę wyszukać specjalistę po nazwisku lub specjalizacji, aby szybko znaleźć odpowiednią osobę na wizytę.
    \item[US2] Jako Użytkownik chcę otrzymać potwierdzenie rezerwacji na e-mail, aby mieć pewność, że proces zakończył się sukcesem.
    \item[US3] Jako Specjalista chcę blokować wybrane dni w grafiku (np. urlop), aby pacjenci nie mogli zapisywać się w czasie mojej nieobecności.
    \item[US4] Jako Specjalista chcę widzieć powód wizyty podany przez pacjenta, aby móc lepiej przygotować się do konsultacji.
    \item[US5] Jako Administrator chcę generować raporty miesięczne dotyczące liczby zrealizowanych wizyt, aby analizować obłożenie systemu.
\end{description}

\section{Acceptance criteria}

\begin{description}[leftmargin=1.5cm, labelwidth=1cm, labelsep=0.5cm]
\textbf{AC1 -- wyszukiwanie specjalistów}
\begin{itemize}[itemsep=2pt]
    \item \textbf{Given} użytkownik znajduje się na stronie wyszukiwania specjalistów
    \item \textbf{When} użytkownik wprowadzi nazwisko lub specjalizację w pole wyszukiwania i zatwierdzi wybór
    \item \textbf{Then} system wyświetla listę specjalistów pasujących do wpisanych kryteriów
\end{itemize}

\vspace{0.5cm}

\textbf{AC2 -- potwierdzenie rezerwacji}
\begin{itemize}[itemsep=2pt]
    \item \textbf{Given} użytkownik pomyślnie przeszedł proces rezerwacji wizyty
    \item \textbf{When} rezerwacja zostanie zatwierdzona w bazie danych
    \item \textbf{Then} system wysyła wiadomość e-mail z potwierdzeniem terminu na adres użytkownika
\end{itemize}

\vspace{0.5cm}

\textbf{AC3 -- blokowanie dni w grafiku}
\begin{itemize}[itemsep=2pt]
    \item \textbf{Given} specjalista jest zalogowany i widzi swój kalendarz pracy
    \item \textbf{When} specjalista zaznaczy wybrany zakres dni jako niedostępny i zapisze zmiany
    \item \textbf{Then} system blokuje możliwość zapisu pacjentów na wizyty w wyznaczonym terminie
\end{itemize}

\vspace{0.5cm}

\textbf{AC4 -- podgląd szczegółów wizyty}
\begin{itemize}[itemsep=2pt]
    \item \textbf{Given} specjalista przegląda szczegóły nadchodzącej wizyty w swoim panelu
    \item \textbf{When} specjalista otworzy kartę konkretnej wizyty
    \item \textbf{Then} system wyświetla pole z opisem powodu wizyty wprowadzonym przez pacjenta
\end{itemize}

\vspace{0.5cm}

\textbf{AC5 -- raportowanie obłożenia}
\begin{itemize}[itemsep=2pt]
    \item \textbf{Given} administrator jest zalogowany w module analitycznym
    \item \textbf{When} administrator wybierze miesiąc i wygeneruje raport obłożenia
    \item \textbf{Then} system tworzy zestawienie zawierające sumaryczną liczbę zrealizowanych wizyt w wybranym okresie
\end{itemize}
\end{description}


\section{Metadane i ocena}

\begin{table}[h]
\begin{tabular}{ll}
\toprule
\textbf{Pole} & \textbf{Wartość} \\
\midrule
Model & \texttt{gemini-3-flash-preview} \\
Tryb & LLM-only \\
Case study & appointment-booking \\
Czas łącznie & 150.4s \\
\midrule
\textit{Ocena (0--3)} & \\
\midrule
Correctness (M1) & 2 \\
Completeness (M2) & 1 \\
Consistency (M3) & 2 \\
Clarity (M4) & 3 \\
Maintainability (M5) & 1 \\
\bottomrule
\end{tabular}
\end{table}


\end{document}